% ============================================================================
%  Subdocumento 11. Planes en operación (D4-153).
% ============================================================================

\begin{resumenApertura}
Este subdocumento fija cómo se mantiene la solución durante los 36 meses de operación sin detener una
flota que rueda a toda hora. Reparte las 640 horas mensuales del equipo de mantención evolutiva entre
corrección, mejoras, deuda técnica y actualización, declara la bolsa anual de 2.688 horas de mantención
evolutiva desde el mes 25, aplica un presupuesto de error de 43,8 minutos al mes a los servicios críticos y ubica cada
prueba de recuperación fuera de los congelamientos del Caso.
\begin{recibe}
  \item Un plan de mantención preventiva, correctiva, evolutiva y normativa con criterios de prioridad y capacidad declarada por tipo.
  \item Una bolsa anual de 2.688 horas de mantención evolutiva desde el mes 25, acumulable, con su procedimiento de solicitud, estimación, aprobación y liquidación.
  \item Una política de presupuesto de error y un análisis sin culpa para cada incidente crítico o alto.
  \item Seis pruebas reales de recuperación y cuatro ejercicios de resiliencia al año, todos fuera de la temporada de fruta y de los demás congelamientos.
\end{recibe}
\end{resumenApertura}

La operación tiene 16 personas en cuatro paquetes del Formulario T-15: la mesa de servicio (13.1, seis
personas), la operación de la plataforma y gestión del servicio (13.2, cuatro), el ciclo de vida del
equipo a bordo (13.3, dos) y la mantención evolutiva (13.4, cuatro). El Subdocumento 10 trata la
atención. Este subdocumento trata lo que se hace sobre la solución misma: la mantención en la sección
11.1, las versiones y la deuda técnica en la 11.2, la confiabilidad en la 11.3, la capacidad y el costo
de la nube en la 11.4 y las pruebas de recuperación en la 11.5.

% ============================================================================
\section{Plan de mantención}\label{sec:11-plan-de-mantencion}
% ============================================================================

Las bases separan cuatro tipos de mantención \transv*{numeral 21.4}{37}. La preventiva sigue un
calendario anual acordado con el mandante, con auditorías técnicas al menos trimestrales e informe. La
correctiva corrige defectos sin costo adicional y dentro de los plazos de resolución del \art{78.2}{40}.
La evolutiva atiende las mejoras que pide el mandante con una bolsa anual de horas. La normativa adecua
la solución a cambios legales sin costo adicional y antes de la vigencia de la norma. La
\tabref{11-capacidad} muestra cómo se reparten las 640 horas mensuales del paquete 13.4 entre esos
trabajos y la deuda técnica.

\begin{tabla}{Capacidad mensual del paquete 13.4 por tipo de trabajo}{11-capacidad}{X R{18mm} R{22mm}}
Tipo de trabajo & Fracción & Horas al mes \\ \midrule
Mantención correctiva y apoyo de nivel 3 & 30 \% & 192 \\
Mantención evolutiva, bolsa del mandante & 35 \% & 224 \\
Reducción de la deuda técnica & 15 \% & 96 \\
Actualización de versiones y mantención normativa & 15 \% & 96 \\
Operación de las innovaciones y desborde de la mesa & 5 \% & 32 \\
Total, cuatro personas de 160 horas & 100 \% & 640 \\
\end{tabla}
\notaTabla{Fuente: elaboración propia sobre el Formulario T-15, paquete 13.4.}

La bolsa evolutiva de 224 horas al mes da 2.688 horas al año, que es el tamaño que declara audIT conforme
a \transv{RT-21.19}{37}. Rige desde el mes 25, por lo que explica el párrafo sobre los meses 21 a 24, y
suma 7.168 horas hasta el mes 56. La composición por perfil es 45 \% de ingeniería de software, 20 \% de
interfaces web y móviles, 20 \% de datos y 15 \% de pruebas, y puede moverse entre perfiles cuando una
solicitud lo pida. Las horas que el mandante no usa en un año pasan completas al siguiente. La tarifa del
excedente va en el Sobre N.º~3, porque el \art{50.2}{29} no admite montos en la oferta técnica. Si la
correctiva de un mes no usa sus 192 horas, el saldo va a la deuda técnica y no a la bolsa, para que la
bolsa del mandante no dependa de cuántos defectos hubo.

Una solicitud evolutiva sigue cinco pasos. La contraparte técnica de Curimón la registra en la mesa como
solicitud de cambio. audIT la estima en horas por perfil dentro de cinco días hábiles. La contraparte la
aprueba, y si pasa de 80 horas la aprueba el Comité de Operación. Se construye con el mismo estándar de
calidad, pruebas y seguridad que el desarrollo original, como exige \transv{RT-21.20}{38}, y pasa por el
mismo pipeline. Cada mes el informe de nivel de servicio liquida las horas usadas por solicitud y el
saldo de la bolsa.

La prioridad sigue un orden fijo. Primero van los incidentes por su severidad del Subdocumento 10.
Después, las vulnerabilidades críticas, que tienen el plazo del \art{21.1}{15}, y las adecuaciones
normativas, ordenadas por la fecha de vigencia de la norma. Las mejoras evolutivas se ordenan por el
efecto en la operación dividido por las horas que cuestan, y la contraparte puede subir una solicitud
cuando la necesita para una temporada.

\paragraph{Los meses 21 a 24} La operación empieza con dos trabajos de la implementación todavía
abiertos. El Formulario T-14 mantiene hasta el mes 24 el retiro del sistema de 2013 (paquete 6.5) y la
migración histórica (paquete 7.2), y los ejecuta el equipo de operación con las horas que ya tienen
asignadas los elementos 6 y 7 de la EDT. La conciliación de capacidad por perfil da para el paquete 13.4
cerca de 372 horas al mes de esos cierres en el mes 21, que bajan a 368 en el 24, más 14 horas de
operación de las innovaciones. El paquete 13.2 pone otras 62 horas al mes, de seguridad y de plataforma.
La \figref{11-transicion} compara cómo se reparte la capacidad del paquete 13.4 en esos cuatro meses y
desde el mes 25.

\begin{figuraNativa}[!ht]{Capacidad mensual del paquete 13.4 en los meses 21 a 24 y desde el mes 25, en horas}{11-transicion}
\begin{tikzpicture}[x=1mm,y=1mm,
  t/.style={font=\sffamily\fontsize{9.5bp}{11bp}\selectfont\color{audit-tinta},anchor=west},
  seg/.style={draw=white,line width=0.6pt}]
  % meses 21 a 24, de abajo hacia arriba
  \fill[seg,fill=audit-marino!75] (0,0) rectangle (24,19.2);
  \fill[seg,fill=audit-gris-claro] (0,19.2) rectangle (24,25.4);
  \fill[seg,fill=audit-tinta!35] (0,25.4) rectangle (24,26.8);
  \fill[seg,fill=audit-turquesa!30] (0,26.8) rectangle (24,64);
  \node[t] at (25.5,9.6) {Correctiva, 192};
  \node[t,align=left] at (25.5,22.6) {Actualización, 62\\ e innovaciones, 14};
  \node[t,align=left] at (25.5,45.4) {Cierre de los\\ paquetes 6.5 y 7.2,\\ 372};
  % desde el mes 25
  \fill[seg,fill=audit-marino!75] (74,0) rectangle (98,19.2);
  \fill[seg,fill=audit-turquesa!70] (74,19.2) rectangle (98,41.6);
  \fill[seg,fill=audit-filete!60] (74,41.6) rectangle (98,51.2);
  \fill[seg,fill=audit-gris-claro] (74,51.2) rectangle (98,60.8);
  \fill[seg,fill=audit-tinta!35] (74,60.8) rectangle (98,64);
  \node[t] at (99.5,9.6) {Correctiva, 192};
  \node[t] at (99.5,30.4) {Bolsa evolutiva, 224};
  \node[t] at (99.5,46.4) {Deuda técnica, 96};
  \node[t] at (99.5,56) {Actualización, 96};
  \node[t] at (99.5,63.2) {Innovaciones y desborde, 32};
  \draw[audit-filete,line width=0.6pt] (-3,0) -- (101,0);
  \node[t,anchor=north] at (12,-1.5) {Meses 21 a 24};
  \node[t,anchor=north] at (86,-1.5) {Meses 25 a 56};
\end{tikzpicture}
\end{figuraNativa}

En los meses 21 a 24 la correctiva conserva sus 192 horas, porque los incidentes no esperan, y las 62
que quedan van a parches y adecuaciones normativas. La bolsa evolutiva y el carril de deuda técnica
empiezan en el mes 25, cuando terminan los cierres. Los perfiles más cargados en esos meses son datos y
calidad: el paquete 13.4 necesita cerca de 146 horas de datos y 77 de pruebas al mes. El Líder de Datos
y la Líder de Calidad aportan 80 y 40 horas desde su dedicación en la operación (Subdocumento 12), y el
resto lo toman los ingenieros del paquete. Las 14 horas de innovaciones de este paquete son la mitad de
las 28 que piden las innovaciones 1 y 2 en operación, 12 del expediente verificable y 16 de la
programación por oportunidad de patio, y la otra mitad la pone el paquete 13.2.

La preventiva tiene dos frentes. En la plataforma, el paquete 13.2 revisa cada mes certificados,
respaldos, alertas y vulnerabilidades, y cada trimestre entrega la auditoría técnica con informe. En el
terreno, el paquete 13.3 revisa los gabinetes y los equipos a bordo cuando pasan por el terminal, porque
el equipo a bordo solo se interviene ahí \caso*{RT-06.01}{32}. El grupo electrógeno de San Bernardo se
prueba cada mes, el clima y la extinción de la sala tienen la mantención de su fabricante y las baterías
de las UPS se reponen en el mes 36. Al salir de servicio, todo medio de almacenamiento se borra con
certificado \transv*{RT-08.17}{20}.

% ============================================================================
\section{Actualización de dependencias y deuda técnica}\label{sec:11-actualizacion-de-dependencias-y-}
% ============================================================================

El Subdocumento 4 declara la versión, el fin de soporte y el plan de actualización de cada producto de
software. La regla de operación es que ningún componente llegue a su fin de soporte sin una actualización
probada: todo componente cuyo soporte termina en los próximos 12 meses entra al plan anual, que es la
ventana de obsolescencia. .NET~10 y PostgreSQL~16, cuyo soporte termina en noviembre de 2028, se
actualizan antes de empezar la operación, dentro del desarrollo de la Etapa 2, como fija el Subdocumento
8. El hardware tiene la misma regla. El iWave G26I declara 10 años de producto, y si la plataforma de
los terminales pierde soporte antes del mes 56, audIT la reemplaza a su cargo \transv*{numeral 1.6}{5}.

Las actualizaciones tienen tres ritmos. Los parches de seguridad se aplican cada mes, o antes si una
vulnerabilidad es crítica. Las versiones menores se agrupan cada trimestre. Las mayores se planifican una
vez al año, con ventana acordada con el mandante y plan de reversión, como exige
\transv{RT-21.22}{38}. Cada artefacto lleva su lista de componentes, y el pipeline la revisa contra las
vulnerabilidades conocidas en cada construcción (Subdocumento 6). Una actualización sigue la promoción
de siempre: desarrollo, QA, preproducción y producción, con la versión anterior lista para volver.

La deuda técnica tiene un registro cuantificado y una fracción fija de la capacidad, como pide
\transv{RT-21.21}{38}. Cada partida del registro dice qué es, qué riesgo trae, cuántas horas cuesta
corregirla y qué cuesta postergarla. El carril de deuda técnica del tablero Kanban que fija el
Subdocumento 6 toma las 96 horas mensuales de la \tabref{11-capacidad}, el 15 \% de la mantención, más
el saldo que deje la correctiva. Esa fracción no viene del 15 \% de los sprints de construcción: sale de
la dotación de operación del Formulario T-15. El registro se revisa cada mes en el Comité de Operación, y
su tamaño total en horas es un indicador del informe mensual.

% ============================================================================
\section{Plan de operación con ingeniería de confiabilidad}\label{sec:11-plan-de-operacion-con-ingenieria}
% ============================================================================

La operación aplica las prácticas de ingeniería de confiabilidad de \textcite{beyer2016}. Cada servicio
tiene un objetivo de nivel de servicio igual al que compromete el Subdocumento 10, y su presupuesto de
error es lo que falta para llegar al 100 \%. Para un servicio crítico de 99,9 \%, el presupuesto es de
43,8 minutos al mes. Si un mes consume más de la mitad, los cambios evolutivos de ese servicio necesitan
una revisión adicional antes de pasar a producción. Si lo agota, se detienen los cambios evolutivos del
servicio, salvo las correcciones y los parches de seguridad, hasta que el mes siguiente cierre dentro del
objetivo. Así el ritmo de cambios se ajusta a la estabilidad medida, y la tasa de cambios fallidos queda
bajo el 5 \% del \art{78.3}{40}.

El trabajo manual repetitivo se mide. Cada mes el paquete 13.2 registra qué parte de su tiempo se fue en
tareas manuales que se repiten y no mejoran el servicio. El límite es la mitad del tiempo
\parencite{beyer2016}, y la meta es bajar de un tercio en el segundo año de operación. Lo que se
automatiza sale de los libros de operación del Subdocumento 10, en el orden que allí se fija.

Todo incidente crítico, y todo alto con efecto en los usuarios, tiene un análisis retrospectivo sin
culpa. El análisis reconstruye la línea de tiempo, busca las causas en el sistema y no en las personas,
y termina en acciones con responsable y plazo. Se entrega al mandante en el plazo del informe de causa
raíz del \art{78.2}{40}, cinco días hábiles para los críticos y diez para los altos. Las acciones se
siguen en el Comité de Operación, y el indicador de reincidencia del \art{78.3}{40}, no más de un evento
por trimestre con la misma causa, mide si funcionaron.

% ============================================================================
\section{Capacidad y optimización en nube}\label{sec:11-capacidad-y-optimizacion-en-nube}
% ============================================================================

El costo de la nube se gestiona con las tres fases del marco FinOps: informar, optimizar y operar
\parencite{finops_framework}. Para informar, todo recurso lleva etiquetas de ambiente, centro de costo y
responsable, con presupuesto y alerta de desviación por ambiente, como exige el \art{16.3}{11}. Cada mes
el informe de operación trae el consumo por ambiente y por servicio contra su presupuesto.

Para optimizar, el tamaño de AKS, PostgreSQL y Redis se ajusta con la prueba de carga de 1,5 veces el
peak antes del hito H5 \transv*{RT-09.06}{21}, y después con el uso real de cada trimestre. Los runners
del pipeline se encienden solo mientras hay trabajos. Las cargas estables, como los nodos de AKS y las
bases de producción, se evalúan para compromisos de capacidad de uno o tres años después de tres meses
de operación sin cambios de tamaño, y la decisión se toma con el mandante, porque la suscripción está a
su nombre. Lo mismo vale para las licencias: los usuarios de GitLab siguen la dotación, 33 en la
primera etapa, 37 en el peak y 16 desde el mes 25, y la capacidad de 100.000 líneas de SonarQube se
revisa cada año contra el código real antes de renovar. En la operación los dos runners del pipeline
corren unas 40 horas al mes cada uno. Para operar, la revisión mensual de costos es parte del Comité de
Operación, y una desviación sobre el presupuesto de un ambiente abre una acción con responsable.

La capacidad se vigila con umbrales que anticipan el límite. IoT Hub agrega una unidad cuando su uso
diario supera el 75 \% de la cuota, Event Hubs cuando el retraso de proceso crece, y PostgreSQL cuando su
almacenamiento pasa del 70 \%. La flota crece a 430 camiones en tres años \caso*{numeral 14.1}{29}, y la
plataforma está dimensionada para tres veces la volumetría inicial \transv*{RT-09.03}{21}, así que el
crecimiento previsto se absorbe con unidades y nodos, sin rediseño.

El hardware crece con la flota en fechas fijas. La flota propia pasa de 148 a 170 camiones, y el
Formulario T-11 prevé para ellos 22 juegos a bordo que compra el mandante: 8 en el mes 24, 7 en el 36 y
7 en el 48. A eso se suman los 19 equipos de reposición del mismo formulario. El paquete 13.3 instala los juegos nuevos en el paso normal de cada camión por el terminal,
igual que en la implantación.

% ============================================================================
\section{Pruebas periódicas de recuperación y resiliencia}\label{sec:11-pruebas-periodicas-de-recuperaci}
% ============================================================================

La recuperación ante desastres se prueba con una conmutación real a Brazil South dos veces al año, con
informe, medición del RPO y del RTO alcanzados y plan de corrección, como exige
\transv{RT-07.07}{17}. La primera prueba se hace antes del paso a producción (Subdocumento 4). Durante
la operación caen en noviembre y en mayo: los meses 22, 28, 34, 40, 46 y 52. El \art{78.3}{40} pide
cumplir el 100 \% del RPO de 15 minutos y del RTO de 4 horas en cada prueba.

Las fechas evitan los congelamientos del \caso{RT-10.05}{32}: la temporada de fruta de diciembre a
abril, Semana Santa, Fiestas Patrias y los nueve días del cierre mensual de liquidaciones. Por eso cada
prueba se hace en la segunda quincena del mes. La \figref{11-calendario} muestra el año tipo de la
operación.

\begin{figuraNativa}[!ht]{Año tipo de la operación: congelamientos del Caso y meses de cada prueba}{11-calendario}
\begin{tikzpicture}[x=8.4mm,y=7mm,
  lab/.style={anchor=east,font=\sffamily\fontsize{9.5bp}{11bp}\selectfont\color{audit-tinta}},
  cel/.style={font=\sffamily\fontsize{9.5bp}{11bp}\selectfont\color{audit-tinta}},
  marca/.style={draw=audit-marino,line width=0.5pt,fill=audit-turquesa!35,minimum width=7mm,minimum height=4.6mm,inner sep=0pt}]
  \fill[audit-gris-claro] (0,0.5) rectangle (4,-8.5);
  \fill[audit-gris-claro] (11,0.5) rectangle (12,-8.5);
  \foreach \m/\c in {1/E,2/F,3/M,4/A,5/M,6/J,7/J,8/A,9/S,10/O,11/N,12/D}{
    \node[cel] at (\m-0.5,0) {\c};
  }
  \draw[audit-filete] (0,-0.5) -- (12,-0.5);
  \node[lab] at (-0.2,0) {Mes};
  \node[lab] at (-0.2,-1) {Congelamiento};
  \node[cel] at (2,-1) {Temporada de fruta};
  \node[cel] at (8.5,-1) {Fiestas Patrias};
  \node[lab] at (-0.2,-2) {Recuperación ante desastres};
  \node[marca] at (4.5,-2) {};
  \node[marca] at (10.5,-2) {};
  \node[lab] at (-0.2,-3) {Ejercicio de resiliencia};
  \node[marca] at (6.5,-3) {};
  \node[marca] at (9.5,-3) {};
  \node[lab] at (-0.2,-4) {Actualización mayor};
  \node[marca,minimum width=24mm] at (6.5,-4) {};
  \node[lab] at (-0.2,-5) {Auditoría trimestral};
  \foreach \m in {3,6,9,12}{ \node[marca,fill=white] at (\m-0.5,-5) {}; }
  \node[lab] at (-0.2,-6) {Prueba de penetración};
  \node[marca,minimum width=15.4mm] at (6,-6) {};
  \node[lab] at (-0.2,-7) {Jornada de actualización};
  \node[marca] at (4.5,-7) {};
  \node[marca] at (9.5,-7) {};
  \node[lab] at (-0.2,-8) {Depósito de fuentes};
  \node[marca,fill=white] at (0.5,-8) {};
  \node[marca,fill=white] at (6.5,-8) {};
\end{tikzpicture}
\end{figuraNativa}

Las pruebas de recuperación caen en mayo y noviembre, los ejercicios de resiliencia en julio y octubre y
la actualización mayor entre junio y agosto, todo fuera de la temporada de fruta. La prueba de
penetración anual de la empresa independiente ocupa junio y julio, y las dos jornadas de actualización
para el personal de Curimón van en mayo y octubre (Subdocumento 10). Las auditorías trimestrales y el
depósito semestral de las fuentes, en enero y julio, son trabajos documentales que no intervienen la
operación, así que pueden caer en cualquier mes. Los cierres del paso Los Libertadores, de junio a septiembre, no cambian este calendario,
porque ninguna prueba depende de un camión en ruta internacional.

Los ejercicios de resiliencia prueban las fallas que no son un desastre regional. En julio se simula la
caída de una zona de disponibilidad de producción y se corta el circuito ExpressRoute para medir la conmutación a la
VPN, que debe ocurrir en hasta un minuto. En octubre un terminal opera 12 horas sin enlace, como exige
\caso{RT-03.10}{31}, y una muestra de equipos a bordo acumula 72 horas sin cobertura y se reconecta a la
vez, para medir que la sincronización no pase de 20 minutos. Además, cada mes se restaura un respaldo
elegido al azar para comprobar que el esquema 3-2-1-1-0 de \transv{RT-07.09}{18} termina en cero errores
de restauración. Cada ejercicio deja un informe con lo medido, lo que falló y el plan de corrección, que
se sigue como cualquier acción de un análisis sin culpa.
